\subsection{Analiza matematyczna dla Algorytmu \ref{alg:jednowymiarowy}}

W niniejszym podrozdziale przeprowadzimy rygorystyczną analizę matematyczną algorytmu zaproponowanego do odzyskiwania funkcji postaci $f(x)=g(a \cdot x)$. Celem tej analizy jest oszacowanie błędu aproksymacji $\|f - \hat{f}\|_\infty$ z odpowiednio wysokim prawdopodobieństwem. Dowód głównego twierdzenia rozbijemy na szereg lematów pomocniczych, które pozwolą nam opanować błędy wynikające z próbkowania skompresowanego, szumu oraz losowego doboru punktów.

\subsubsection{Główne twierdzenie o błędzie aproksymacji}
Zacznijmy od sformułowania głównego rezultatu, którego dowód będzie zwieńczeniem tej części pracy.

\begin{tw}[Zbieżność Algorytmu \ref{alg:jednowymiarowy} \cite{fornasier2012learning}]\label{tw:zbieznoscalgorytmu1}
Niech $0 < s < 1$ oraz $\log d \le m_\Phi \le [\log 6]^{-2}d$. Wówczas istnieje stała $c_1'$, taka że wykorzystując $m_\mathcal{X} \cdot (m_\Phi + 1)$ ewaluacji funkcji $f$, Algorytm \ref{alg:jednowymiarowy} definiuje funkcję $\hat{f}: B_{\mathbb{R}^d}(1+\bar{\epsilon}) \to \mathbb{R}$, która z prawdopodobieństwem co najmniej
$$1 - \left(e^{-c_1' m_\Phi} + e^{-\sqrt{m_\Phi d}} + 2e^{-\frac{2m_\mathcal{X} s^2 \alpha^2}{C_2^4}}\right)$$
spełnia nierówność
$$\|f - \hat{f}\|_\infty \le 2C_2(1+\bar{\epsilon})\frac{\nu_1}{\sqrt{\alpha(1-s)} - \nu_1},$$
gdzie
$$\nu_1 = C' \left( \left[ \frac{m_\Phi}{\log(d/m_\Phi)} \right]^{1/2 - 1/q} + \frac{\epsilon}{\sqrt{m_\Phi}} \right)$$
oraz $C'$ zależy tylko od stałych $C_1$ i $C_2$ definiujących model.
\end{tw}

Aby udowodnić to twierdzenie, wprowadzimy teraz niezbędne fakty z teorii próbkowania skompresowanego oraz rachunku prawdopodobieństwa.

\subsubsection{Narzędzia z teorii próbkowania skompresowanego}

Fundamentem działania algorytmu jest zastosowanie macierzy pomiarowej spełniającej własność ograniczonej izometrii (RIP). Własność tę opisuje poniższe twierdzenie, stanowiące specjalizację twierdzenia 1.2 z \cite{wojtaszczyk2012l1}, dla macierzy Bernoulliego.

\begin{tw}\label{tw:aproksymacjaK-czlonowa}
    Załóżmy, że $\Phi$ jest losową macierzą wymiaru $m \times d$, w której wszystkie wpisy są niezależnymi zmiennymi o rozkładzie Bernoulliego przeskalowanymi przez $1/\sqrt{m}$.
    \begin{itemize}
        \item Dla $0 < \delta < 1$ istnieją stałe $c_1, c_2 > 0$ takie, że z prawdopodobieństwem co najmniej $1 - e^{-c_1 m}$, macierz $\Phi$ posiada własność RIP (Restricted Isometry Property):
            $$(1-\delta)\|x\|_{l_2^d}^2 \le \|\Phi x\|_{l_2^m}^2 \le (1+\delta)\|x\|_{l_2^d}^2,$$ 
        dla wszystkich $x \in \mathbb{R}^d$ takich, że $\#\text{supp}(x) \le c_2 m / \log(d/m)$.
        \item Jeśli $d > [\log 6]^2 m$, to istnieją stałe $C, c_1', c_2' > 0$ takie, że z prawdopodobieństwem co najmniej $1 - e^{-c_1' m} - e^{-\sqrt{md}}$ dla każdego wektora $x \in \mathbb{R}^d$, szumu $\epsilon \in \mathbb{R}^m$ i naturalnego $K \le c_2' m / \log(d/m)$ zachodzi ograniczenie odzyskiwania sygnału:
            $$\|\Delta(\Phi x + \epsilon) - x\|_{l_2^d} \le C(K^{-1/2}\sigma_K(x)_{l_1^d} + \max\{\|\epsilon\|_{l_2^m}, \sqrt{\log d}\|\epsilon\|_{l_\infty^m}\}),$$ 
        gdzie $\sigma_K(x)_{l_1^d}$ jest błędem najlepszej $K$-członowej aproksymacji wektora $x$ w normie $l_1$.
    \end{itemize}
\end{tw}

Twierdzenie to podajemy bez dowodu. Jego pierwsza część jest powszechnie znanym wynikiem z teorii próbkowania skompresowanego \cite{baraniuk2008simple, fornasier2010numerical} , natomiast druga część stanowi relatywnie nowy rezultat, wynikający z prac Wojtaszczyka \cite{wojtaszczyk2012l1} połączonych z ustaleniami DeVore'a, Petrovej i Wojtaszczyka \cite{devore2009instance}.

Aplikując Twierdzenie \ref{tw:aproksymacjaK-czlonowa} do układu równań z Algorytmu \ref{alg:jednowymiarowy}, uzyskujemy kluczowy wniosek szacujący błąd odzyskania kolumn $x_j$:

\begin{wn}\label{wn:bladodzyskaniakolumn}
Niech $d > [\log 6]^2 m_\Phi$.
\begin{itemize}
    \item Z prawdopodobieństwem co najmniej $1 - (e^{-c_1' m_\Phi} + e^{-\sqrt{m_\Phi d}})$, wszystkie wektory $\hat{x}_j = \Delta(y_j)$ wyliczone w Algorytmie \ref{alg:jednowymiarowy} spełniają:
$$\|x_j - \hat{x}_j\|_{l_2^d} \le C\left(\left[\frac{m_\Phi}{\log(d/m_\Phi)}\right]^{1/2 - 1/q} + \max\{\|\epsilon_j\|_{l_2^{m_\Phi}}, \sqrt{\log d}\|\epsilon_j\|_{l_\infty^{m_\Phi}}\}\right).$$ 
    \item Jeśli dodatkowo $m_\Phi \ge \log d$, to z tym samym prawdopodobieństwem zachodzi:
$$\|x_j - \hat{x}_j\|_{l_2^d} \le C'\left(\left[\frac{m_\Phi}{\log(d/m_\Phi)}\right]^{1/2 - 1/q} + \frac{\epsilon}{\sqrt{m_\Phi}}\right)$$ 
gdzie stałe zależą tylko od $C_1$ i $C_2$.
\end{itemize}
\end{wn}

\begin{proof}
Twierdzenie \ref{tw:aproksymacjaK-czlonowa} aplikujemy do równania $y_j = \Phi x_j + \epsilon_j$ dla $K \le c_2' m_\Phi / \log(d/m_\Phi)$. Ze względu na strukturę $X$, każda jej kolumna to przeskalowany wektor $a^T$, tj. $x_j = g'(a \cdot \xi_j)a^T$. Błąd najlepszej aproksymacji szacujemy więc następująco:
$$K^{-1/2}\sigma_K(x_j)_{l_1^d} \le |g'(a \cdot \xi_j)| \cdot \|a\|_{l_q^d} \cdot K^{1/2 - 1/q} \le C_1 C_2 \left[\frac{m_\Phi}{\log(d/m_\Phi)}\right]^{1/2 - 1/q},$$ 
co kończy dowód części pierwszej Wniosku \ref{wn:bladodzyskaniakolumn}.
Aby udowodnić drugą część tego wniosku, szacujemy rozmiar błędów Taylora $\epsilon_j$:
$$\|\epsilon_j\|_{l_\infty^{m_\Phi}} \le \frac{\epsilon \|g''\|_\infty}{2m_\Phi} \left(\sum_{k=1}^d |a_k|\right)^2 \le \frac{C_1^2 C_2}{2m_\Phi} \epsilon,$$ 
$$\|\epsilon_j\|_{l_2^{m_\Phi}} \le \sqrt{m_\Phi} \|\epsilon_j\|_{l_\infty^{m_\Phi}} \le \frac{C_1^2 C_2}{2\sqrt{m_\Phi}} \epsilon.$$ 
Zatem maksimum z norm wynosi:
$$\max\{\|\epsilon_j\|_{l_2^{m_\Phi}}, \sqrt{\log d}\|\epsilon_j\|_{l_\infty^{m_\Phi}}\} \le \frac{C_1^2 C_2}{2\sqrt{m_\Phi}}\epsilon \cdot \max\left\{1, \sqrt{\frac{\log d}{m_\Phi}}\right\}.$$ 
Przy założeniu $m_\Phi \ge \log d$, czynnik z maksimum znika (ogranicza się przez 1), co daje tezę Wniosku \ref{wn:bladodzyskaniakolumn}. 
\end{proof}


\subsubsection{Stabilność kierunkowa i koncentracja miary}

Znamy już błąd bezwzględny wyznaczenia $x_j$. Teraz potrzebujemy pokazać, że błąd ten nie zniekształca zbytnio samej orientacji (kierunku) wyznaczonego wektora.

\begin{lm}[Stabilność podprzestrzeni – przypadek jednowymiarowy \cite{fornasier2012learning}]\label{lm:stabilnoscpodprzestrzenijednowymiarowy}
Ustalmy $\hat{x} \in \mathbb{R}^d$, wektor znormalizowany $a \in \mathbb{S}^{d-1}$, skalary $\gamma \neq 0$ oraz wektor szumu $n \in \mathbb{R}^d$ o normie $\|n\|_{l_2^d} \le \nu_1 < |\gamma|$. Jeśli przyjmiemy, że $\hat{x} = \gamma a + n$, to:
$$\left\| \text{sign}(\gamma)\frac{\hat{x}}{\|\hat{x}\|_{l_2^d}} - a \right\|_{l_2^d} \le \frac{2\nu_1}{\|\hat{x}\|_{l_2^d}}.$$ 
\end{lm}

\begin{proof}
Wielokrotnie wykorzystując nierówność trójkąta (oraz jej postać odwrotną) i pamiętając, że $a \in \mathbb{S}^{d-1}$, otrzymujemy:
$$\left\| \text{sign}(\gamma)\frac{\hat{x}}{\|\hat{x}\|_{l_2^d}} - a \right\|_{l_2^d} \le \left\| \text{sign}(\gamma)\frac{\hat{x}}{\|\hat{x}\|_{l_2^d}} - \frac{|\gamma|a}{\|\hat{x}\|_{l_2^d}} \right\|_{l_2^d} + \left\| \frac{|\gamma|a}{\|\hat{x}\|_{l_2^d}} - a \right\|_{l_2^d}$$ 
$$\le \frac{\|n\|_{l_2^d}}{\|\hat{x}\|_{l_2^d}} + \left| \frac{|\gamma|}{\|\hat{x}\|_{l_2^d}} - 1 \right| \le \frac{2\nu_1}{\|\hat{x}\|_{l_2^d}}.$$ 
co stanowi dowód lematu.
\end{proof}

Ponieważ oszacowanie w Lemacie \ref{lm:stabilnoscpodprzestrzenijednowymiarowy} jest tym lepsze, im większa jest norma $\|\hat{x}_j\|_{l_2^d}$, algorytm słusznie szuka indeksu maksymalnego $j_0$. Aby zagwarantować, że to maksimum jest wystarczająco duże i nie "utonie" w szumie (tzn. $\nu_1 < |\gamma|$), potrzebujemy użyć koncentracji probabilistycznej:

\begin{tw}[Nierówność Hoeffdinga \cite{ledoux2001concentration}]\label{tw:nierownoschoeffdinga}
Niech $X_1, \dots, X_m$ będą niezależnymi zmiennymi losowymi. Załóżmy, że zmienne te są prawie na pewno ograniczone, tzn. $\mathbb{P}\{X_j - \mathbb{E}X_j \in [a_j, b_j]\} = 1$. Wówczas:
$$\mathbb{P}\left\{\left|\sum_{j=1}^m X_j - \mathbb{E}\left(\sum_{j=1}^m X_j\right)\right| \ge t\right\} \le 2e^{-\frac{2t^2}{\sum_{j=1}^m (b_j - a_j)^2}}.$$ 
\end{tw}

Zastosowanie tego faktu prowadzi nas do kluczowego lematu o wielkości gradientu.

\begin{lm}[\cite{fornasier2012learning}]\label{lm:wielkoscgradientu}
Ustalmy $0 < s < 1$. Wówczas z prawdopodobieństwem $1 - 2e^{-\frac{2m_\mathcal{X} s^2 \alpha^2}{C_2^4}}$ zachodzi nierówność:
$$\max_{j=1,\dots,m_\mathcal{X}} |g'(a \cdot \xi_j)| \ge \sqrt{\alpha(1-s)},$$ 
gdzie $\alpha := \mathbb{E}_\xi(|g'(a \cdot \xi_j)|^2)$.
\end{lm}

\begin{proof}
Z definicji stałych z modelu (równania (8) i (9)) mamy $\mathbb{E}X_j = \int_{\mathbb{S}^{d-1}} |g'(a \cdot \xi)|^2 d\mu_{\mathbb{S}^{d-1}}(\xi) \ge \alpha > 0$  oraz ograniczenie nośnika $X_j - \mathbb{E}X_j \in [-\alpha, C_2^2 - \alpha]$. Aplikując do tych zmiennych nierówność Hoeffdinga otrzymujemy:
$$\mathbb{P}\left\{\left| \sum_{j=1}^{m_\mathcal{X}} |g'(a \cdot \xi_j)|^2 - m_\mathcal{X} \alpha \right| \ge s m_\mathcal{X} \alpha\right\} \le 2e^{-\frac{2m_\mathcal{X} s^2 \alpha^2}{C_2^4}}.$$ 
Zatem, ze wskazanym prawdopodobieństwem zachodzi:
$$\frac{1}{m_\mathcal{X}} \sum_{j=1}^{m_\mathcal{X}} |g'(a \cdot \xi_j)|^2 \ge \alpha(1-s).$$ 
Gdyby dla każdego indeksu $j$ zachodziło $|g'(a \cdot \xi_j)|^2 < \alpha(1-s)$, powyższa średnia nie mogłaby być osiągnięta. Wynika stąd bezpośrednio, że relacja ta musi zachodzić dla maksimum, co kończy dowód. 
\end{proof}

\subsubsection{Dowód Twierdzenia \ref{tw:zbieznoscalgorytmu1}}

Mając już wszystkie potrzebne do tego narzędzia, możemy przejść do dowodu Twierdzenia \ref{tw:zbieznoscalgorytmu1}.

\begin{proof}
Na mocy Lematu \ref{lm:wielkoscgradientu} wiemy, że $|g'(a \cdot \xi_{j_0})| \ge \sqrt{\alpha(1-s)}$ z prawdopodobieństwem $1 - 2e^{-\frac{2m_\mathcal{X} s^2 \alpha^2}{C_2^4}}$. Dodatkowo, Wniosek \ref{wn:bladodzyskaniakolumn} w połączeniu z Lematem \ref{lm:stabilnoscpodprzestrzenijednowymiarowy} (podstawiając szum i błędy algorytmu) gwarantuje nam, że odzyskany kierunek $\hat{a}$ spełnia ograniczenie:
$$\|\text{sign}(g'(a \cdot \xi_{j_0}))\hat{a} - a\|_{l_2^d} \le \frac{2\nu_1}{\sqrt{\alpha(1-s)} - \nu_1},$$ 
z łącznym prawdopodobieństwem $1 - \left(e^{-c_1' m_\Phi} + e^{-\sqrt{m_\Phi d}} + 2e^{-\frac{2m_\mathcal{X} s^2 \alpha^2}{C_2^4}}\right)$.

Mając to oszacowanie podprzestrzeni, możemy wyznaczyć globalny błąd aproksymacji funkcji $f$. Dla dowolnego punktu $x \in B_{\mathbb{R}^d}(1+\bar{\epsilon})$:
$$|f(x) - \hat{f}(x)| = |g(a \cdot x) - \hat{g}(\hat{a} \cdot x)| = |g(a \cdot x) - f(\hat{a}^T \hat{a} \cdot x)| = |g(a \cdot x) - g(a \cdot \hat{a}^T \hat{a} \cdot x)|.$$ 
Wykorzystując ograniczenie gładkościowe funkcji ($C_2$) oraz fakt, że $\hat{a}(I_d - \hat{a}^T \hat{a}) = 0$, kontynuujemy ciąg nierówności, wprowadzając nieznany znak:
$$\le C_2 |a \cdot x - a \cdot [\hat{a}^T \hat{a}] \cdot x| = C_2 |a \cdot (I_d - \hat{a}^T \hat{a})x|$$ 
$$= C_2 |(a - \text{sign}(g'(a \cdot \xi_{j_0}))\hat{a}) \cdot (I_d - \hat{a}^T \hat{a})x|$$ 
Stosując nierówność Cauchy'ego-Schwarza i zauważając, że projekcja $I_d - \hat{a}^T \hat{a}$ ma normę co najwyżej 1:
$$\le C_2 \|a - \text{sign}(g'(a \cdot \xi_{j_0}))\hat{a}\|_{l_2^d} \cdot \|x\|_{l_2^d}.$$ 
Ponieważ wektor $x$ operuje wewnątrz zdefiniowanej powłoki $1+\bar{\epsilon}$, jego norma nie przekracza tej wartości. Ostatecznie otrzymujemy więc:
$$|f(x) - \hat{f}(x)| \le 2C_2(1+\bar{\epsilon})\frac{\nu_1}{\sqrt{\alpha(1-s)} - \nu_1},$$ 
co w pełni udowadnia twierdzenie.
\end{proof}

\subsubsection{Uwagi o złożoności informacyjnej}
Na koniec tej analizy warto udowodnić, że wyznaczona tolerancja szumu $\alpha$ nie zapada się wykładniczo, unikając w ten sposób klątwy wymiarowości. Opiera się to na twierdzeniu określającym postać analityczną rzutu miary $\mu_k$ \cite{rudin1980function}  oraz na rozwinięciu funkcji z uogólnionego modelu Taylora.

\begin{tw}[Rzut miary sferycznej \cite{rudin1980function}]\label{tw:rzutmiary}
Niech $1 \le k < d$ będą liczbami naturalnymi. Wówczas miara $\mu_k$ na kuli jednostkowej $B_{\mathbb{R}^k}$, będąca rzutem (push-forward) miary $\mu_{\mathbb{S}^{d-1}}$, wyraża się wzorem:
$$d\mu_k(y) = \frac{\Gamma(d/2)}{\pi^{k/2}\Gamma((d-k)/2)}(1 - \|y\|_{l_2^k}^2)^{\frac{d-2-k}{2}}dy.$$ 
\end{tw}

Ograniczenie to pozwala nam przejść do fundamentalnego sformułowania dotyczącego redukcji błędu.

\begin{stw}[Spadek wielomianowy parametru $\alpha$ \cite{fornasier2012learning}]\label{stw:spadekwielomianowy}
Załóżmy, że $g \in C^{M+2}$ w otoczeniu zera oraz $\frac{d^l}{dx^l}g(0) = 0$ dla $l=1,\dots,M$. Wówczas parametr koncentracji energii $\alpha(d)$ wykazuje spadek wielomianowy:
$$\alpha(d) = \frac{\Gamma(d/2)}{\pi^{1/2}\Gamma((d-1)/2)} \int_{-1}^1 |g'(y)|^2 (1 - y^2)^{\frac{d-3}{2}}dy = \mathcal{O}(d^{-M})$$ 
przy $d \to \infty$.
\end{stw}

\begin{proof}
Korzystając z rozwinięcia Taylora wokół zera dla funkcji $g'$, wszystkie wyrazy aż do rzędu $M$ redukują się, co daje $|g'(y)|^2 = \left(\frac{1}{M!}\frac{d^{M+1}}{dx^{M+1}}g(0)\right)^2 y^{2M} + \mathcal{O}(y^{2M+1})$. Biorąc pod uwagę asymetryczność kuli i własności funkcji Gamma dla całek z momentami, odrzucamy nieparzyste składniki błędu. Stosując do pozostałej całki aproksymację Stirlinga ($\Gamma(z) \approx \sqrt{2\pi/z}(z/e)^z$) , wyrażenie dla rzutu miary przy $d \to \infty$ transformuje się do asymptoty:
$$\alpha(d) = \mathcal{O}\left( \frac{\Gamma(d/2)\Gamma((1+2M)/2)}{\Gamma((d+2M)/2)} \right) \approx d^{-M},$$ 
co dowodzi potęgowej (a nie wykładniczej) degradacji użytecznego sygnału. 
\end{proof}